% !TeX root = RJwrapper.tex
\title{fastsae: Frequentist and Bayesian Small Area Estimation in R}


\author{by Ridson Al Farizal P. and Azka Ubaidillah}

\maketitle

\abstract{%
Small area estimation (SAE) combines survey data with auxiliary information to estimate parameters of domains whose samples are too small for direct estimation. We present fastsae, an R package that brings the main area-level and unit-level models of official statistics together in one interface. The frequentist part fits the Fay-Herriot model, its spatial and spatio-temporal extensions, the Battese-Harter-Fuller unit-level model and the two-fold sub-area level model, with analytical and bootstrap mean squared errors, in compiled code that never forms a dense inverse of the marginal covariance matrix. The hierarchical Bayes part fits area-level models with Gaussian, beta, gamma, binomial, Poisson and negative binomial likelihoods, and unit-level and two-fold models with Gaussian, binomial and Poisson likelihoods, with spatial, temporal and space-time random effects, by integrated nested Laplace approximations or, through the same call, by Hamiltonian Monte Carlo in JAX. For beta, gamma and negative binomial data the dispersion of each area is fixed by the design-based sampling variance of its direct estimate, and the default priors were chosen by simulation and are documented. Further functions benchmark estimates to published totals at one or two levels, classify them by relative standard error, compare models, and produce maps and spreadsheet exports. We describe the methods and their implementation, check the estimates against the packages sae and emdi, compare runtimes with sae, emdi and tipsae, report a design-based simulation, and illustrate a complete workflow with poverty rates of the health districts of Emilia-Romagna.
}

\section{Introduction}\label{sec:intro}

Official statistics are increasingly asked for indicators at levels of geography or population detail that national surveys were not designed for: poverty rates of districts, unemployment of municipalities, or health indicators of subregions \citep{rao2015small, pfeffermann2013new}. Direct estimates from such domains rest on few observations and have large sampling variances. Small area estimation (SAE) reduces these variances by linking the domains through a model with auxiliary information, at the area level when only domain aggregates are available \citep{fay1979estimates} and at the unit level when auxiliary variables are known for the sampled units and as population means \citep{battese1988error}. The basic models have been extended to spatially correlated area effects \citep{petrucci2006small, pratesi2008small}, to repeated surveys \citep{marhuenda2013small}, and, for proportions, rates and counts, to non-Gaussian likelihoods, either through empirical best predictors of generalized linear mixed models \citep{boubeta2016empirical} or by hierarchical Bayes (HB) estimation, which also accounts for the uncertainty about the variance parameters \citep{rao2015small, janicki2020properties, fabrizi2011hierarchical}.

Producing small area statistics routinely raises three practical requirements. First, the methods must scale to the thousands of areas, and the repeated fits of bootstrap mean squared error (MSE) estimation, of national production systems. Second, the HB models that suit bounded and count indicators should not require each analyst to write and diagnose an MCMC sampler. Third, published small area estimates are usually benchmarked to direct estimates of larger domains and accompanied by indicators of their reliability, so these steps belong in the same workflow as the model fit.

Several R packages implement parts of this (Table \ref{tab:pkg-comparison}). \CRANpkg{sae} \citep{molina2015sae} covers the frequentist area-level models, including spatial and spatio-temporal Fay-Herriot models, and the unit-level model in R code. \CRANpkg{emdi} \citep{kreutzmann2019emdi} adds transformations, the empirical best predictor for poverty indicators, benchmarking to a national total, maps and spreadsheet export. \CRANpkg{hbsae} \citep{boonstra2022hbsae} fits Gaussian area- and unit-level HB models by numerical integration over the variance parameter and benchmarks estimates to any set of linear constraints. \CRANpkg{tipsae} \citep{denicolo2024tipsae} fits beta-type area-level models with spatial and temporal effects in \pkg{Stan}, \CRANpkg{mcmcsae} \citep{boonstra2022mcmcsae} fits a wide class of multilevel models by MCMC, and \CRANpkg{SUMMER} \citep{li2019changes, li2025summer} fits area- and cluster-level models with INLA, focusing on demographic and health surveys.

\begin{table}[!h]
\centering
\caption{\label{tab:pkg-comparison}Small area estimation packages for R. EBLUP: empirical best linear unbiased predictor; EB(P): empirical best (predictor); HB: hierarchical Bayes; HMC: Hamiltonian Monte Carlo. Taken from the package documentation (sae 1.3, emdi 2.2.3, hbsae 1.2, tipsae 1.0.3, mcmcsae 0.8.2, SUMMER 2.0.0).}
\centering
\fontsize{7}{9}\selectfont
\begin{tabular}[t]{>{\raggedright\arraybackslash}p{2.2cm}>{\raggedright\arraybackslash}p{1.75cm}>{\raggedright\arraybackslash}p{1.75cm}>{\raggedright\arraybackslash}p{1.75cm}>{\raggedright\arraybackslash}p{1.75cm}>{\raggedright\arraybackslash}p{1.75cm}>{\raggedright\arraybackslash}p{1.75cm}>{\raggedright\arraybackslash}p{1.75cm}}
\toprule
Feature & fastsae & sae & emdi & hbsae & tipsae & mcmcsae & SUMMER\\
\midrule
Inference & REML/ML, HB (INLA, HMC) & REML/ML & REML/ML & REML, HB (num. integration) & HB (Stan HMC) & HB (Gibbs, HMC) & HB (INLA)\\
Area-level model & yes & yes & yes & yes & beta-type & yes & yes\\
Spatial / spatio-temporal EBLUP & SAR / SAR + AR(1) & SAR / SAR + AR(1) & SAR / -- & -- & -- & -- & --\\
Unit-level model & EBLUP, HB & EBLUP, EB & EBLUP, EBP & EBLUP, HB & -- & HB & HB\\
HB likelihoods & Gaussian, beta, gamma, binomial, Poisson, neg. binomial & -- & -- & Gaussian & beta, flexible beta, 0/1-inflated beta, ext. beta & Gaussian, binomial, multinomial, Poisson, neg. binomial, gamma & Gaussian, binomial, beta-binomial\\
HB spatial effects & ICAR, BYM, BYM2, Leroux, SAR, SLM & -- & -- & -- & ICAR & ICAR, BYM2, Leroux & ICAR, BYM2\\
Benchmarking & ratio, difference, MSE-weighted, logit; one or two levels & -- & raking, ratio, MSE-adjusted; one total & linear constraints & raking, ratio, double & -- & national totals\\
Maps / export & yes / spreadsheet & -- / -- & yes / spreadsheet & -- / -- & yes / CSV & -- / -- & yes / --\\
Compiled code & C++ (OpenMP) & -- & -- & -- & Stan & C++ & --\\
\bottomrule
\end{tabular}
\end{table}

\CRANpkg{fastsae} covers these models and steps in one interface and was written with the first two requirements in mind. Its contributions are:

\begin{enumerate}
\def\labelenumi{\arabic{enumi}.}
\tightlist
\item
  Compiled implementations of the Fay-Herriot, spatial, spatio-temporal, unit-level and two-fold models with analytical and bootstrap MSEs that never form a dense inverse of the marginal covariance matrix, and run the bootstrap in parallel with results that do not depend on the number of threads.
\item
  One HB interface, \texttt{hb\_area()}, \texttt{hb\_unit()} and \texttt{hb\_twofold()}, for six likelihoods at the area level and three at the unit and sub-area levels, with spatial, temporal and space-time effects, fitted by INLA \citep{rue2009approximate} or by the No-U-Turn sampler of NumPyro \citep{hoffman2014no, phan2019composable} through interchangeable engine objects. For beta, gamma and negative binomial data the dispersion of each area is fixed by the sampling variance of its direct estimate, and the default priors are documented together with the simulations that led to them.
\item
  Post-estimation functions for benchmarking at one or two levels, reliability classification, model comparison, maps and export, which accept the output of every model.
\end{enumerate}

The remainder of the paper describes the \hyperref[sec:methods]{methods} and their \hyperref[sec:design]{implementation}, reports \hyperref[sec:simulations]{benchmarks and a design-based simulation}, and illustrates a \hyperref[sec:applications]{complete workflow}.

\section{Methods}\label{sec:methods}

Throughout, areas are indexed by \(d = 1, \dots, D\), \(\boldsymbol{x}_d\) is a \(p\)-vector of area-level covariates, and \(\psi_d\) is the design-based sampling variance of the direct estimate \(y_d\) of the area parameter \(\theta_d\).

\subsection{Frequentist models}\label{sec:methods-freq}

The Fay-Herriot model \citep{fay1979estimates} is
\begin{equation}
y_d = \boldsymbol{x}_d^\top \boldsymbol{\beta} + u_d + e_d, \quad u_d \overset{\mathrm{iid}}{\sim} \mathcal{N}(0, \sigma_u^2), \quad e_d \overset{\mathrm{ind}}{\sim} \mathcal{N}(0, \psi_d),
\label{eq:fh}
\end{equation}
with \(\psi_d\) known. Its empirical best linear unbiased predictor (EBLUP) is \(\hat{\theta}_d = \hat{\gamma}_d y_d + (1 - \hat{\gamma}_d)\boldsymbol{x}_d^\top\hat{\boldsymbol{\beta}}\) with \(\hat{\gamma}_d = \hat{\sigma}_u^2 / (\hat{\sigma}_u^2 + \psi_d)\), where \(\hat{\boldsymbol{\beta}}\) is the generalized least squares estimator at the REML or ML estimate \(\hat{\sigma}_u^2\). The MSE is estimated by the second-order approximation \(g_{1d} + g_{2d} + 2g_{3d}\), which \citet{prasad1990estimation} derived for a moment estimator of \(\sigma_u^2\) and \citet{datta2000unified} showed to hold for REML; for ML, \citet{datta2000unified} add a bias term. The (RE)ML estimate is truncated at zero, as in \CRANpkg{sae}. With few areas or large sampling variances this happens often; the EBLUPs then equal the synthetic estimates \(\boldsymbol{x}_d^\top\hat{\boldsymbol{\beta}}\) and the MSE estimates can understate their error \citep{li2010adjusted}. \texttt{eblup\_fh()} reports such fits with a message; the adjusted likelihoods proposed for this problem \citep{li2010adjusted, yoshimori2014new} are not implemented. For positive indicators, \texttt{eblup\_fh(transform\ =\ "log")} fits the model to \(\log y_d\) and back-transforms with the conditional-mean correction \(\exp(\hat{\mu}_d + \hat{g}_{1d}/2)\) and a delta-method MSE, which omits the higher-order terms of \citet{slud2006mean}.

The spatial Fay-Herriot model replaces the independent area effects by a simultaneous autoregressive (SAR) process, \(\boldsymbol{u} = \rho\boldsymbol{W}\boldsymbol{u} + \boldsymbol{\varepsilon}\) with \(\boldsymbol{\varepsilon} \sim \mathcal{N}(\boldsymbol{0}, \sigma_u^2\boldsymbol{I})\) and a row-standardized proximity matrix \(\boldsymbol{W}\) \citep{petrucci2006small, pratesi2008small}. The marginal covariance is then \(\boldsymbol{V} = \boldsymbol{G} + \boldsymbol{\Psi}\) with \(\boldsymbol{G} = \sigma_u^2[(\boldsymbol{I} - \rho\boldsymbol{W})^\top(\boldsymbol{I} - \rho\boldsymbol{W})]^{-1}\) and \(\boldsymbol{\Psi} = \mathrm{diag}(\psi_d)\), and the EBLUP is \(\hat{\boldsymbol{\theta}} = \boldsymbol{X}\hat{\boldsymbol{\beta}} + \hat{\boldsymbol{G}}\hat{\boldsymbol{V}}^{-1}(\boldsymbol{y} - \boldsymbol{X}\hat{\boldsymbol{\beta}})\). The process is defined on all areas of \(\boldsymbol{W}\), so areas without a sample are predicted as well. \((\sigma_u^2, \rho)\) are estimated by REML or ML Fisher scoring on \(|\rho| \le 0.999\), with a projected iteration when the unconstrained one ends at a negative variance. The analytical MSE is the second-order approximation \(g_{1d} + g_{2d} + 2g_{3d} - g_{4d}\) of \citet{singh2005spatio}, where \(g_{4d}\) accounts for the nonlinearity of \(\boldsymbol{V}\) in \(\rho\), and the parametric and non-parametric bootstrap MSEs are those of \citet{molina2009bootstrap}, as in \CRANpkg{sae} \citep{molina2015sae}. Two choices depart from \CRANpkg{sae}. When \(\hat{\rho}\) lies on the boundary of its range the expansion is not valid, so \(\rho\) is then treated as fixed in the MSE; and for ML fits, \(g_{3d}\) and the bias term use the ML rather than the REML information, because the latter made the MSE explode for \(\hat{\sigma}_u^2 \approx 0\) and \(\hat{\rho} \approx 1\) in simulations. For areas without a sample, the MSE of the predictor includes \(g_{3d}\) and the bias of the plug-in \(g_{1d}\).

For repeated surveys, \(t = 1, \dots, T\), the spatio-temporal model of \citet{marhuenda2013small} is
\begin{equation}
y_{dt} = \boldsymbol{x}_{dt}^\top\boldsymbol{\beta} + u_{1d} + u_{2dt} + e_{dt},
\label{eq:stfh}
\end{equation}
where \(\boldsymbol{u}_1 = (u_{11}, \dots, u_{1D})^\top\) follows the SAR process of the spatial model with parameters \((\sigma_1^2, \rho_1)\), the effects of each area follow a stationary AR(1) process, \(u_{2dt} = \rho_2 u_{2d,t-1} + \varepsilon_{2dt}\) with \(\varepsilon_{2dt} \overset{\mathrm{iid}}{\sim} \mathcal{N}(0, \sigma_2^2)\), independently across areas, and \(e_{dt} \overset{\mathrm{ind}}{\sim} \mathcal{N}(0, \psi_{dt})\). With \texttt{model\ =\ "S"} the \(u_{2dt}\) are independent. As in \citet{marhuenda2013small} and \CRANpkg{sae}, the variance parameters are estimated by REML only, and the MSE is estimated analytically by \(g_{1dt} + g_{2dt} + 2g_{3dt} - g_{4dt}\) or by a parametric bootstrap that re-estimates the model in every replicate. The panel need not be complete: the model is defined on the full \(D \times T\) grid, the variance parameters are estimated from the observed cells, and missing periods and areas that were never sampled are predicted from the observed ones \citep{saei2005out}, with analytical and bootstrap MSEs for every cell.

At the unit level, the nested error model of \citet{battese1988error}, \(y_{dj} = \boldsymbol{x}_{dj}^\top\boldsymbol{\beta} + u_d + e_{dj}\) for unit \(j\) of area \(d\), with \(u_d \overset{\mathrm{iid}}{\sim} \mathcal{N}(0, \sigma_u^2)\) and \(e_{dj} \overset{\mathrm{iid}}{\sim} \mathcal{N}(0, \sigma_e^2)\), gives the EBLUP of the area mean \(\hat{\bar{Y}}_d = f_d\bar{y}_d + (1 - f_d)\{\bar{\boldsymbol{X}}_{d,r}^\top\hat{\boldsymbol{\beta}} + \hat{\gamma}_d(\bar{y}_d - \bar{\boldsymbol{x}}_d^\top\hat{\boldsymbol{\beta}})\}\), where \(f_d = n_d/N_d\) is the sampling fraction, \(\bar{\boldsymbol{X}}_{d,r}\) the covariate mean of the non-sampled units and, here, \(\hat{\gamma}_d = \hat{\sigma}_u^2/(\hat{\sigma}_u^2 + \hat{\sigma}_e^2/n_d)\). Its MSE is estimated by the parametric bootstrap of \citet{gonzalez2008bootstrap}.

When the areas are divided into sub-areas, such as districts within provinces, with direct estimates \(y_{dj}\) of the sub-area parameters \(\theta_{dj}\), \texttt{eblup\_twofold()} fits the sub-area level model of \citet{torabi2014small},
\begin{equation}
y_{dj} = \theta_{dj} + e_{dj}, \qquad \theta_{dj} = \boldsymbol{x}_{dj}^\top\boldsymbol{\beta} + v_d + u_{dj},
\label{eq:twofold}
\end{equation}
with area effects \(v_d \overset{\mathrm{iid}}{\sim} \mathcal{N}(0, \sigma_v^2)\), sub-area effects \(u_{dj} \overset{\mathrm{iid}}{\sim} \mathcal{N}(0, \sigma_u^2)\) and \(e_{dj} \overset{\mathrm{ind}}{\sim} \mathcal{N}(0, \psi_{dj})\). A sub-area without a sample borrows from the sampled sub-areas of its area through \(\hat{v}_d\). The variance components are estimated by REML or ML, the MSE analytically (\(g_1 + g_2 + 2g_3\), with the bias term for ML) or by parametric bootstrap, and the function also returns the area means \(\theta_d = \sum_j w_{dj}\theta_{dj}\) with an MSE that includes the covariances between the sub-area prediction errors.

\subsection{Hierarchical Bayes models}\label{sec:methods-hb}

\subsubsection{Latent Gaussian models}\label{latent-gaussian-models}

The HB functions formulate each model as a latent Gaussian model \citep{rue2009approximate}. For \texttt{hb\_area()},
\begin{equation}
y_d \mid \mu_d \sim \pi(y_d \mid \mu_d, \phi_d), \qquad g(\mu_d) = \boldsymbol{x}_d^\top\boldsymbol{\beta} + u_d,
\label{eq:lgm}
\end{equation}
where \(g\) is the identity, logit or log link of the family, \(\phi_d\) a dispersion parameter (below), and \(\boldsymbol{u}\) a zero-mean Gaussian Markov random field: independent effects or one of the spatial structures below, plus temporal and space-time terms for panel data. \texttt{hb\_unit()} applies the same construction to unit records with an area effect, and \texttt{hb\_twofold()} to sub-areas within areas. Areas without a direct estimate stay in the model with a missing response, and their estimates are the posterior summaries of \(\mu_d\), which borrow from the covariates and, with a spatial or temporal effect, from the neighbouring areas or periods. The INLA engine approximates the posterior marginals by nested Laplace approximations and numerical integration over the hyperparameters; \CRANpkg{fastsae} uses the full Laplace approximation of the latent marginals by default instead of INLA's simplified one \citep{rue2009approximate}. Gaussian models, and the Leroux and spatial lag effects for every family, are fitted in INLA's classic mode \citep{rue2009approximate}, the other models in its default compact mode, which leaves the linear predictors out of the latent field \citep{vanniekerk2023new}. These choices rest on comparisons with exact posteriors (numerical integration over the variance parameters) and with long \pkg{Stan} runs of the same models and priors: for Gaussian models the compact mode overstated the posterior standard deviations by up to 18\%, whereas in the classic mode the posterior means were within 0.07 posterior standard deviations of the reference and the standard deviations within 0.91 to 1.03 times; for binomial, Poisson, beta and gamma models both modes matched \pkg{Stan} and the compact mode was much faster. The scripts are in \texttt{benchmarks/validation/} of the package repository. The JAX engine samples the same model with the No-U-Turn sampler.

\subsubsection{Dispersion from the sampling variances}\label{sec:methods-dispersion}

When the sampling variances \(\psi_d\) are known, \CRANpkg{fastsae} fixes the dispersion of each area so that the model variance of \(y_d\) equals \(\psi_d\):

\begin{itemize}
\tightlist
\item
  \textbf{Gaussian}: the observation variance is \(\psi_d\), which gives the Fay-Herriot model \eqref{eq:fh}.
\item
  \textbf{Beta}: for proportions with a logit link, \(\mathrm{Var}(y_d) = \mu_d(1 - \mu_d)/(1 + \phi_d)\), which gives the precision of \citet{janicki2020properties},
  \begin{equation}
  \phi_d = \max\left\{\frac{\mu_d(1 - \mu_d)}{\psi_d} - 1, \, 1\right\}.
  \label{eq:betaphi}
  \end{equation}
  It depends on the unknown \(\mu_d\). Replacing \(\mu_d\) by \(y_d\) makes the model variance \(\psi_d\mu_d(1 - \mu_d)/\{y_d(1 - y_d)\}\), too small wherever \(y_d\) is far from \(\mu_d\); in a simulation with proportions near 0.08 the 95\% intervals then covered 81--84\% of the true values. \CRANpkg{fastsae} therefore fits the model twice and evaluates \eqref{eq:betaphi} at the posterior means of the first fit, which gave 90--94\%. \CRANpkg{tipsae} instead makes \(\phi_d\) a function of \(\mu_d\) inside the model \citep{denicolo2024tipsae}; this covered 93--95\% in the same simulation, but makes the likelihood bimodal, with a mode of low precision near zero, and 11 of 240 fits did not converge (\(\hat{R} > 1.05\)). With known sample sizes instead of variances, \(\phi_d = n_d - 1\), which assumes simple random sampling; for complex designs, effective sample sizes \(n_d/\mathrm{deff}_d\) should be supplied instead.
\item
  \textbf{Gamma}: for positive skewed indicators with a log link, \(\mathrm{Var}(y_d) = \mu_d^2/s_d\) gives the shape \(s_d = y_d^2/\psi_d\), the inverse squared coefficient of variation of the direct estimate. The shape is evaluated at the direct estimate, so the model variance is \(\psi_d\mu_d^2/y_d^2\) rather than \(\psi_d\), the same approximation that the two-stage precision avoids for the beta model.
\item
  \textbf{Negative binomial}: \(\mathrm{Var}(y_d) = \mu_d + \mu_d^2/k_d\) gives the size \(k_d = \hat{\mu}_d^2/(\psi_d - \hat{\mu}_d)\), again at the posterior means of a first fit; areas with \(\psi_d \le \hat{\mu}_d\), i.e.~without overdispersion relative to the Poisson model, get a near-Poisson size.
\item
  \textbf{Binomial and Poisson}: the response is a count, the numbers of trials are sample sizes and population exposures enter as an offset. Direct estimates from a weighted sample are not counts of a simple random sample; for them, the effective number of trials \(\tilde{n}_d = \hat{p}_d(1 - \hat{p}_d)/\psi_d\) and \(\tilde{n}_d\hat{p}_d\) successes, which reproduce the design-based variance, are the usual inputs \citep{chen2014use} (\texttt{hb\_area()} accepts proportions with \texttt{trials} and rounds the counts), or the beta model with \(\psi_d\) can be used.
\end{itemize}

Without sampling variances, one dispersion parameter is estimated. With one direct estimate per area it trades off against the area effect, so its prior determines the intervals (Table \ref{tab:priors}), and \texttt{hb\_area()} warns.

\subsubsection{Spatial and space-time effects}\label{spatial-and-space-time-effects}

The area effect supports the intrinsic conditional autoregressive (ICAR) model of \citet{besag1991bayesian}, the BYM model with independent and ICAR components, the scaled BYM2 model \(u_d = \sigma_u(\sqrt{1 - \phi}\,v_d + \sqrt{\phi}\,u^*_d)\) of \citet{riebler2016intuitive}, where \(u^*\) is a scaled ICAR field, the Leroux model with precision \(\tau[(1 - \lambda)\boldsymbol{I} + \lambda(\mathrm{diag}(\boldsymbol{W}\boldsymbol{1}) - \boldsymbol{W})]\) \citep{leroux2000estimation}, and SAR models for the area effect, as in the spatial Fay-Herriot model, or as a spatial lag model. Temporal effects can be random walks of order one or two, AR(1) or independent, and space-time interactions can be area-specific trends, a spatial field with AR(1) dynamics, or one of the four types of \citet{knorrheld2000bayesian}.

\subsubsection{Default priors}\label{sec:methods-priors}

With few areas, or one observation per area, the data say little about the variance and dispersion parameters, and their priors decide how strongly the estimates are shrunk and how wide the intervals are. Table \ref{tab:priors} lists the defaults, which both engines share.

\begin{table}[!h]
\centering
\caption{\label{tab:priors}Default priors of hb\_area(), hb\_unit() and hb\_twofold(). All intrinsic Gaussian Markov random fields are scaled, so that $\sigma$ is a typical marginal standard deviation for every structure. Every prior can be changed through inla\_engine().}
\centering
\begin{tabular}[t]{>{\raggedright\arraybackslash}p{6cm}>{\raggedright\arraybackslash}p{8cm}}
\toprule
Parameter & Prior\\
\midrule
Coefficients $\boldsymbol{\beta}$ & intercept flat; others $\mathcal{N}(0, 1000)$\\
Standard deviations of random effects (area, sub-area, time, space-time; $\sigma_e$ of hb\_unit()) & half-normal$(s)$, $s = \mathrm{sd}(\boldsymbol{y})$ for Gaussian models and $s = 1$ on the link scale otherwise\\
BYM2 mixing parameter $\phi$ & PC prior with $P(\phi < 0.5) = 0.5$\\
Leroux $\lambda$; SAR and SLM $\rho$ & uniform on the admissible range\\
AR(1) correlation $\rho$ & $\mathcal{N}(0, 1/0.15)$ on $\log\{(1 + \rho)/(1 - \rho)\}$\\
Gamma shape, negative binomial size (no $\psi_d$) & half-normal$(10)$\\
Beta precision (no $\psi_d$, no $n_d$) & exponential with rate 0.1\\
Gaussian residual standard deviation (no $\psi_d$) & half-normal$(s)$\\
\bottomrule
\end{tabular}
\end{table}

The half-normal prior is a weakly informative choice for hierarchical standard deviations \citep{gelman2006prior}. Its scale is tied to the response, because a penalized complexity (PC) prior \(P(\sigma > U) = \alpha\) \citep{simpson2017penalising} with the common \(U = 1\) depends on the unit of \(y\): in a Fay-Herriot simulation with 30 areas the 95\% intervals under PC\((1, 0.01)\) covered 94\%, 85\% and 38\% of the true means for the same data in three units. The scaled half-normal prior covered as well as the EBLUP intervals (93--97\% for 15 to 60 areas, against 93--96\%), whereas PC priors with \(U = \mathrm{sd}(\boldsymbol{y})\) covered 85--94\%, because they put much mass near \(\sigma = 0\) \citep{sorbye2014scaling}. On the logit and log scales, half-normal\((1)\) covered at least as well as the half-normal\((2.5)\) default of \CRANpkg{tipsae} and estimated the area standard deviation with less bias. For the dispersion parameters estimated without \(\psi_d\), the INLA defaults put the gamma shape near 80 and the negative binomial size near 110--180 whatever the truth, and intervals for overdispersed data covered 52--92\% of the true means; the half-normal\((10)\) prior covered 90--99\%, at the cost of too wide intervals when the sampling noise is small, which is why \(\psi_d\) should be supplied whenever it is known. These simulations and their code are documented in \texttt{vignette("hb-priors")}.

\subsection{Benchmarking}\label{sec:methods-bench}

Small area estimates \(\hat{\theta}_d\) are often required to agree with direct estimates of larger domains \citep{rao2015small, stefan2021benchmarked}. Let areas \(d\) belong to groups \(g\) (provinces, say) with weights \(w_d\) (population shares) normalized within groups, \(\tilde{w}_d = w_d / \sum_{d' \in g} w_{d'}\), and let \(T_g\) be the group targets. \texttt{benchmark\_sae()} adjusts the estimates so that \(\sum_{d \in g}\tilde{w}_d\hat{\theta}^{\mathrm{B}}_d = T_g\) by

\begin{itemize}
\tightlist
\item
  ratio adjustment, \(\hat{\theta}^{\mathrm{B}}_d = \hat{\theta}_d T_g / \sum_{d' \in g}\tilde{w}_{d'}\hat{\theta}_{d'}\) \citep{rao2015small};
\item
  difference adjustment, \(\hat{\theta}^{\mathrm{B}}_d = \hat{\theta}_d + T_g - \sum_{d' \in g}\tilde{w}_{d'}\hat{\theta}_{d'}\);
\item
  MSE-weighted adjustment, which minimizes \(\sum_d(\hat{\theta}^{\mathrm{B}}_d - \hat{\theta}_d)^2/\mathrm{MSE}(\hat{\theta}_d)\) under the constraint and so moves imprecise estimates most \citep{datta2011bayesian, bell2013benchmarking}:
  \begin{equation}
  \hat{\theta}^{\mathrm{B}}_d = \hat{\theta}_d + \frac{\tilde{w}_d\,\mathrm{MSE}(\hat{\theta}_d)}{\sum_{d' \in g}\tilde{w}_{d'}^2\,\mathrm{MSE}(\hat{\theta}_{d'})}\Big(T_g - \sum_{d' \in g}\tilde{w}_{d'}\hat{\theta}_{d'}\Big);
  \label{eq:bench-opt}
  \end{equation}
\item
  a shift \(\alpha_g\) on the logit scale, \(\sum_{d \in g}\tilde{w}_d\,\mathrm{logit}^{-1}\{\mathrm{logit}(\hat{\theta}_d) + \alpha_g\} = T_g\), solved by one-dimensional root finding, which keeps proportions inside \((0, 1)\).
\end{itemize}

With a national target \(T\) as well, the group targets are first adjusted so that their weighted mean equals \(T\), and the areas are then benchmarked to the adjusted group targets. \texttt{benchmark\_sae()} also returns a first-order approximation of the MSE of the benchmarked estimates, which uses the covariances between the area prediction errors of the fitted model. The area-level and two-fold EBLUP functions and \texttt{hb\_area()} can also benchmark within the fit (argument \texttt{self\_benchmark}): the EBLUP functions with the augmented model of \citet{wang2008small}, and \texttt{hb\_area()} by adjusting every posterior draw, so that the intervals refer to the benchmarked estimates.

\subsection{Reliability classification}\label{sec:methods-reliability}

Statistical offices publish small area estimates by tiers of their relative standard error, \(\mathrm{RSE}_d = \sqrt{\mathrm{MSE}(\hat{\theta}_d)}/|\hat{\theta}_d|\). \texttt{add\_reliability\_flags()} labels an estimate reliable for \(\mathrm{RSE}_d \le c_1\), to be used with caution for \(c_1 < \mathrm{RSE}_d \le c_2\), and unreliable above \(c_2\). The defaults \(c_1 = 25\%\) and \(c_2 = 50\%\) are those of BPS-Statistics Indonesia; other cut-offs can be supplied. For HB fits the posterior standard deviation takes the place of \(\sqrt{\mathrm{MSE}(\hat{\theta}_d)}\); the two measure different things, which \texttt{compare\_sae()} points out when it compares an HB with an EBLUP fit.

\section{Implementation}\label{sec:design}

Table \ref{tab:functions} lists the main functions. All model functions return objects of class \texttt{"fastsae"}, with the estimates, sampling variances, MSEs or posterior summaries and relative standard errors of all areas in one data frame, sampled or not, so the post-estimation functions accept every model. \texttt{diagnose()} computes the goodness-of-fit and bias diagnostics of \citet{brown2001evaluation} and residual checks, and \texttt{step\_sae()} selects covariates by stepwise search, with the information criteria of \citet{marhuenda2014information} for EBLUP models and the logarithmic conditional predictive ordinate, WAIC or DIC for HB models. The frequentist functions call C++ code {[}\CRANpkg{Rcpp}, \citet{eddelbuettel2011rcpp}; \CRANpkg{RcppArmadillo}, \citet{eddelbuettel2014rcpparmadillo}; \CRANpkg{RcppEigen}{]} that never forms a dense inverse of the marginal covariance matrix. The spatial model works with sparse Cholesky factors of \((\boldsymbol{I} - \rho\boldsymbol{W})^\top(\boldsymbol{I} - \rho\boldsymbol{W})\) and of \((\boldsymbol{I} - \rho\boldsymbol{W})^\top(\boldsymbol{I} - \rho\boldsymbol{W}) + \sigma_u^2\boldsymbol{\Psi}^{-1}\) through the Woodbury identity; the spatio-temporal and two-fold models compute REML scores, Fisher information and the analytical MSE from \(T \times T\), \(D \times D\) and low-rank blocks; and the unit-level model works with area-level sufficient statistics, profiling \(\boldsymbol{\beta}\) and \(\sigma_e^2\) out of the (restricted) likelihood. Bootstrap replicates run in parallel with OpenMP; their random numbers are drawn beforehand from R's generator, so the results depend on the seed but not on the number of threads.

The HB functions translate formulas, adjacency information and sampling variances into a model and pass it to an engine object. \texttt{inla\_engine()} (the default) holds the INLA strategy, the priors and further arguments of \texttt{INLA::inla()}; \texttt{jax\_engine()} holds the device (CPU or GPU), the number of chains and draws and the sampler settings, and runs NumPyro through \CRANpkg{reticulate}. Both engines return the same object, so \texttt{print()}, \texttt{summary()}, \texttt{plot()}, \texttt{coef()}, \texttt{fitted()} and the post-estimation functions do not depend on how a model was estimated. INLA is not on CRAN and Python is optional, so both are suggested dependencies, as are \CRANpkg{sf} for maps and a spreadsheet writer for export; the frequentist and post-estimation functions need neither INLA nor Python.

\begin{table}[!h]
\centering
\caption{\label{tab:functions}Main functions of fastsae.}
\centering
\begin{tabular}[t]{>{\raggedright\arraybackslash}p{5.5cm}>{\raggedright\arraybackslash}p{8.5cm}}
\toprule
Function & Purpose\\
\midrule
eblup\_fh(), eblup\_sfh(), eblup\_stfh() & Fay-Herriot, spatial and spatio-temporal EBLUP with analytical or bootstrap MSE\\
eblup\_bhf(), eblup\_twofold() & Unit-level and two-fold EBLUP\\
hb\_area() & Area-level HB models: six likelihoods, spatial, temporal and space-time effects\\
hb\_unit(), hb\_twofold() & Unit-level and two-fold HB models\\
inla\_engine(), jax\_engine() & Estimation engines for the HB functions\\
benchmark\_sae() & Ratio, difference, MSE-weighted and logit benchmarking at one or two levels\\
add\_reliability\_flags() & Reliability tiers by relative standard error\\
compare\_sae(), step\_sae() & Agreement and ranking of two models; covariate selection\\
diagnose(), autoplot() & Model diagnostics and plots\\
map\_sae(), export\_sae() & Choropleth maps; spreadsheet export\\
\bottomrule
\end{tabular}
\end{table}

\section{Benchmarks and simulations}\label{sec:simulations}

All timings were measured on an Apple M4 (10 cores, 16 GB) under R 4.6.1 with R's reference BLAS; OpenMP used its default number of threads, which matters only for the bootstrap, not timed here. The scripts are in the \texttt{benchmarks/} directory of the package repository and store their results in \texttt{inst/extdata/}, from which this article computes its tables and figures.

\subsection{Agreement with reference implementations}\label{sec:agreement}

\begin{table}[!h]
\centering
\caption{\label{tab:agreement}Largest relative difference, $\max|a - b| / \max|b|$, between fastsae ($a$) and sae 1.3 or emdi 2.2.3 ($b$) fitted to the same data: regression coefficients, variance parameters (including $\rho$), EBLUPs and analytical MSEs. sae has no analytical MSE for the spatio-temporal model, and both packages estimate the MSE of the unit-level model by bootstrap.}
\centering
\fontsize{8}{10}\selectfont
\begin{tabular}[t]{llllrrrr}
\toprule
Model & Data & Method & Reference & Coefficients & Variances & EBLUP & MSE\\
\midrule
Fay-Herriot & mys & REML & sae & 3e-14 & 6e-15 & 5e-15 & 6e-15\\
Fay-Herriot & mys & REML & emdi & 2e-07 & 5e-06 & 5e-07 & 3e-06\\
Fay-Herriot & mys & ML & sae & 9e-15 & 4e-15 & 2e-15 & 4e-15\\
Fay-Herriot & mys & ML & emdi & 3e-07 & 6e-06 & 6e-07 & 4e-06\\
Spatial Fay-Herriot & simulated, D = 100 & REML & sae & 7e-16 & 1e-15 & 6e-16 & 5e-15\\
Spatial Fay-Herriot & simulated, D = 100 & REML & emdi & 6e-16 & 5e-16 & 6e-16 & 5e-15\\
Spatial Fay-Herriot & simulated, D = 100 & ML & sae & 2e-15 & 3e-14 & 1e-14 & 2e-02\\
Spatial Fay-Herriot & simulated, D = 100 & ML & emdi & 5e-15 & 1e-13 & 4e-14 & 2e-02\\
Spatio-temporal Fay-Herriot & simulated, D = 100, T = 5 & REML & sae & 6e-16 & 9e-15 & 2e-15 & --\\
Battese-Harter-Fuller & cornsoybean & REML & sae & 1e-09 & 7e-09 & 1e-09 & --\\
\bottomrule
\end{tabular}
\end{table}

Table \ref{tab:agreement} compares \CRANpkg{fastsae} with \CRANpkg{sae} and \CRANpkg{emdi} on the same data: the sampled areas of the example data \texttt{mys}, simulated spatial and spatio-temporal data with \(D = 100\) areas (the design of the timings below), and the \texttt{cornsoybean} unit-level data of \citet{battese1988error}. Against \CRANpkg{sae}, all estimates and the analytical MSEs of REML fits agree to rounding error; the unit-level fit agrees to the tolerance of the optimizer. The differences from \CRANpkg{emdi}, of order \(10^{-6}\), reflect its convergence tolerance. The ML MSE of the spatial model differs by 1.8\% because \CRANpkg{fastsae} uses the ML information in \(g_{3d}\) and the bias term (see \hyperref[sec:methods-freq]{Frequentist models}).

\subsection{Frequentist models}\label{frequentist-models}

\begin{table}[!h]
\centering
\caption{\label{tab:freq-table}Median time and increase of the peak resident memory of the R process for point estimates and analytical MSEs, except for the spatio-temporal model ($T = 5$ periods), where only point estimates are timed because sae has no analytical MSE for it. All packages fitted the same data sets. Memory increases of a few tens of MB vary from run to run. $^\dagger$The fit did not converge within the iteration limit.}
\centering
\begin{tabular}[t]{rrrrrrr}
\toprule
\multicolumn{1}{c}{ } & \multicolumn{3}{c}{Time (s)} & \multicolumn{3}{c}{Peak memory (MB)} \\
\cmidrule(l{3pt}r{3pt}){2-4} \cmidrule(l{3pt}r{3pt}){5-7}
$D$ & fastsae & sae & emdi & fastsae  & sae  & emdi \\
\midrule
\addlinespace[0.3em]
\multicolumn{7}{l}{\textbf{Fay-Herriot}}\\
\hspace{1em}30 & 0.00033 & 0.00084 & 0.0037 & 21.89 & 44.12 & 116\\
\hspace{1em}100 & 0.00037 & 0.0032 & 0.02 & 22.91 & 57.19 & 106\\
\hspace{1em}500 & 0.0005 & 0.14 & 1.63 & 21.53 & 30.09 & 112\\
\hspace{1em}1,000 & 0.00064 & 0.80 & 13.47 & 1.91 & 39.02 & 222\\
\addlinespace[0.3em]
\multicolumn{7}{l}{\textbf{Spatial Fay-Herriot}}\\
\hspace{1em}30 & 0.0031 & 0.01 & 0.01 & 77.28 & 174 & 178\\
\hspace{1em}100 & 0.0073 & 0.05 & 0.06 & 49.73 & 66.70 & 74.31\\
\hspace{1em}500 & 0.11 & 3.08 & 3.45 & 57.08 & 98.80 & 128\\
\hspace{1em}1,000 & 0.52 & 24.10 & 25.21 & 78.67 & 239 & 310\\
\addlinespace[0.3em]
\multicolumn{7}{l}{\textbf{Spatio-temporal Fay-Herriot, point estimates}}\\
\hspace{1em}30 & 0.07 & 0.70$^\dagger$ & -- & 44.02 & 127 & --\\
\hspace{1em}100 & 0.11 & 1.95 & -- & 52.97 & 137 & --\\
\hspace{1em}500 & 5.21 & 124 & -- & 116 & 1,114 & --\\
\hspace{1em}1,000 & 39.84 & 1,156 & -- & 495 & 3,684 & --\\
\bottomrule
\end{tabular}
\end{table}

We timed the three area-level models on simulated data with five covariates and \(D\) up to 1,000 areas (Table \ref{tab:freq-table}), with the same data sets for all packages. Each combination of model, package and \(D\) ran in its own R process: times are medians from the \CRANpkg{bench} package, and memory is the increase of the peak resident memory of the process over that of a process that had loaded the same packages and data, so it includes memory allocated by compiled code. At \(D = 1{,}000\), \CRANpkg{fastsae} fitted the Fay-Herriot model in 0.00064 s against 0.80 s for \CRANpkg{sae} and 13.47 s for \CRANpkg{emdi}, and the spatial model in 0.52 s against 24.10 s and 25.21 s. For the spatio-temporal model, \CRANpkg{sae} provides point estimates and a bootstrap MSE but no analytical MSE, so we timed point estimates: 39.84 s for \CRANpkg{fastsae} and 19 minutes for \CRANpkg{sae}; with the analytical MSE, \CRANpkg{fastsae} took about as long (39.34 s). The Fisher scoring of \CRANpkg{sae} did not converge within its 100 iterations for \(D = 30, 50\), so these times are those of 100 iterations, whereas \CRANpkg{fastsae} switches to a constrained iteration when the classical one fails. With 30 areas every fit took less than a second, so the differences matter mainly when a model is fitted many times, as in simulation studies.

\subsection{Beta models: fastsae and tipsae}\label{sec:bench-beta}

\begin{table}[!h]
\centering
\caption{\label{tab:beta-table}Median time in seconds of beta models with sampling variances, fitted by fastsae (INLA, two-stage precision) and tipsae (Stan, 4 chains of 2,000 iterations). The spatio-temporal model has $T = 5$ periods and was run up to $D = 50$.}
\centering
\begin{tabular}[t]{rrrrrrr}
\toprule
\multicolumn{1}{c}{ } & \multicolumn{2}{c}{Independent} & \multicolumn{2}{c}{Spatial (ICAR)} & \multicolumn{2}{c}{Spatio-temporal} \\
\cmidrule(l{3pt}r{3pt}){2-3} \cmidrule(l{3pt}r{3pt}){4-5} \cmidrule(l{3pt}r{3pt}){6-7}
$D$ & fastsae & tipsae & fastsae & tipsae & fastsae & tipsae\\
\midrule
30 & 2.85 & 2.31 & 2.66 & 48.71 & 2.93 & 574\\
50 & 2.56 & 4.35 & 2.51 & 55.15 & 3.80 & 716\\
100 & 2.73 & 10.94 & 2.67 & 89.36 & -- & --\\
250 & 3.27 & 33.58 & 2.75 & 145 & -- & --\\
500 & 3.31 & 82.75 & 3.03 & 232 & -- & --\\
1,000 & 3.73 & 180 & 4.28 & 362 & -- & --\\
\bottomrule
\end{tabular}
\end{table}

For beta models we compared \texttt{hb\_area(family\ =\ "beta")} with \CRANpkg{tipsae} \citep{denicolo2024tipsae} on data simulated by \texttt{sim\_area\_data()}, with sampling variances \(\psi_d = y_d(1 - y_d)/26\) computed from the direct estimates as a survey would, and \CRANpkg{tipsae} with \texttt{type\_disp\ =\ "var"} and its default sampler settings (Table \ref{tab:beta-table}). The two packages do not fit the same model: their priors differ (see \hyperref[sec:methods-priors]{Default priors}), and the precision is fixed at a first-stage estimate in \CRANpkg{fastsae} but a function of the model mean in \CRANpkg{tipsae}. The comparison is therefore between what a user obtains from each package with its defaults. At \(D = 1{,}000\) the independent model took 3.73 s in \CRANpkg{fastsae} and 180 s in \CRANpkg{tipsae}, and the spatial model 4.28 s and 362 s. At these sizes most of an INLA fit is fixed cost, such as starting the INLA program, and the two-stage precision fits the model twice, so the runtime of \CRANpkg{fastsae} grows slowly with \(D\).

To compare the estimates, we fitted both packages to the \(D = 1{,}000\) data set of Table \ref{tab:beta-table} and compared them with the true area means (Figure \ref{fig:concordance}). The posterior means of the two packages correlated at \(r = 0.9988\) with a mean absolute difference of 0.0096, and the posterior standard deviations of \CRANpkg{fastsae} were 0.99 times those of \CRANpkg{tipsae} in the median area (90\% of the areas between 0.88 and 1.17). Against the true means, the RMSE was 0.0678 for \CRANpkg{fastsae} and 0.0684 for \CRANpkg{tipsae}, and the 95\% intervals covered 93.3\% and 91.1\% of the true means, with mean widths 0.252 and 0.259.

\begin{figure}

{\centering \includegraphics[alt={Two scatter plots. Left: posterior means from fastsae on the x axis against tipsae on the y axis, between 0 and 1, close to the identity line. Right: posterior standard deviations from fastsae against tipsae.}]{fastsae_files/figure-latex/concordance-1} 

}

\caption{Posterior means (left) and posterior standard deviations (right) of a beta model with $D = 1{,}000$ areas fitted by fastsae and tipsae. The dashed lines are the identity.}\label{fig:concordance}
\end{figure}

\subsection{Estimation engines}\label{sec:bench-engines}

\begin{table}[!h]
\centering
\caption{\label{tab:engines}Time in seconds of a second (warm) fit of hb\_area() with an ICAR area effect, by estimation engine. JAX: NumPyro NUTS, 2 chains of 1,000 draws after 500 warm-up iterations; tipsae: Stan defaults (4 chains of 2,000 iterations), beta model only. Times are for these run lengths, not per effective sample.}
\centering
\begin{tabular}[t]{rrrrr}
\toprule
$D$ & INLA & JAX (CPU) & JAX (Apple GPU) & tipsae\\
\midrule
\addlinespace[0.3em]
\multicolumn{5}{l}{\textbf{Gaussian}}\\
\hspace{1em}100 & 0.73 & 4.51 & 136 & --\\
\hspace{1em}500 & 3.12 & 32.48 & 245 & --\\
\hspace{1em}1,000 & 0.80 & 92.72 & 377 & --\\
\hspace{1em}2,500 & 1.41 & 841 & 725 & --\\
\addlinespace[0.3em]
\multicolumn{5}{l}{\textbf{Beta}}\\
\hspace{1em}100 & 2.60 & 24.00 & 221 & 66.73\\
\hspace{1em}500 & 3.11 & 146 & 281 & 152\\
\hspace{1em}1,000 & 3.72 & 374 & 376 & 235\\
\hspace{1em}2,500 & 6.32 & 1,269 & 819 & 620\\
\bottomrule
\end{tabular}
\end{table}

The same \texttt{hb\_area()} model can be estimated by INLA or by the No-U-Turn sampler of NumPyro on JAX, on the CPU or a GPU. Table \ref{tab:engines} compares the engines on Gaussian and beta models with an ICAR area effect on a regular grid with queen adjacency; for the beta model, \CRANpkg{tipsae} was run with its ICAR effect. Each engine ran in its own R process and each fit was run twice on the same data. At \(D = 2,500\) INLA fitted the spatial beta model in 6.32 s, against 620 s for \CRANpkg{tipsae} and 1,269 s for JAX on the CPU. The JAX engine is a way to check INLA by MCMC rather than a faster alternative to it. Its time grows faster with \(D\) than that of INLA, because it represents the ICAR effect through the eigenvectors of the graph Laplacian, which costs \(O(D^2)\) per gradient, and for the beta model it runs the sampler twice, once for each stage of the precision. On the CPU it took 0.4 times the time of \CRANpkg{tipsae} at \(D = 100\) and 2.0 times at \(D = 2,500\). The Apple GPU has a fixed cost of about 2 minutes per sampler run, as each step of the tree-building loop of NUTS synchronizes GPU and CPU; it pays off only for the largest models. Running NumPyro on Apple GPUs required replacing two operations that \pkg{jax-metal} cannot compile (\texttt{popcnt} and \texttt{lax.cond} with constant outputs). The posterior means of JAX correlated with those of INLA at \(r \ge 0.9999\) with median ratios of posterior standard deviations between 1.00 and 1.02. The largest split \(\hat{R}\) of the JAX fits was 1.016, slightly above the threshold of 1.01 recommended by \citet{vehtari2021rank}, so longer chains than these two of 1,000 draws are advisable for inference. Not all \CRANpkg{tipsae} fits converged: the largest \(\hat{R}\) was 3.71, at \(D = 100\), in line with the bimodal likelihood described under \hyperref[sec:methods-dispersion]{Dispersion from the sampling variances}, so its times refer to the run length and not to a converged posterior.

\subsection{Design-based simulation}\label{sec:mc}

To assess the estimators under repeated sampling from a finite population, we generated a population of 250,824 individuals in \(D = 500\) areas on a \(20 \times 25\) grid, with a binary outcome that is logistic in two covariates with ICAR and independent area effects, and drew \(R = 200\) stratified simple random samples of 20 to 50 units per area. The direct estimates were the sample proportions, with a continuity adjustment that keeps them inside \((0, 1)\), and their variances under simple random sampling with the finite population correction. Each sample was fitted by the Fay-Herriot and spatial Fay-Herriot EBLUPs with analytical MSEs (\texttt{eblup\_fh()}, and \texttt{eblup\_sfh()} with the row-standardized contiguity matrix), and by beta models with an ICAR area effect and the sampling variances, by \texttt{hb\_area()} with its defaults and by \CRANpkg{tipsae} (\texttt{type\_disp\ =\ "var"}, 2 chains of 1,000 iterations after 500 warm-up iterations). The intervals were \(\hat{\theta}_d \pm 1.96\sqrt{\mathrm{mse}_d}\) for the EBLUPs and the 95\% credible intervals for the HB fits. The script is \texttt{benchmarks/sim\_mc\_finite\_pop\_d500.R}.

\begin{table}[!h]
\centering
\caption{\label{tab:mc}Design-based simulation ($D = 500$, $R = 200$): root mean squared error, relative RMSE and mean absolute relative bias, averaged over areas; coverage of the 95\% intervals (Monte Carlo standard error across replicates in parentheses) and their mean width; median over areas of the ratio of the mean estimated MSE (posterior variance for the HB fits) to the empirical MSE; and time per fit with 5 replicates running in parallel.}
\centering
\fontsize{9}{11}\selectfont
\begin{tabular}[t]{llrrrrrrr}
\toprule
  & Estimator & RMSE & RRMSE (\%) & MARB (\%) & Coverage (\%) & Width & Variance ratio & Time (s)\\
\midrule
Direct & Direct & 0.0725 & 27.4 & 21.6 & 94.8 (0.1) & 0.292 & 1.09 & --\\
FH & Fay-Herriot EBLUP & 0.0518 & 20.0 & 16.6 & 94.4 (0.1) & 0.207 & 1.27 & 0.0096\\
SFH & Spatial Fay-Herriot EBLUP & 0.0497 & 19.1 & 15.9 & 94.2 (0.1) & 0.198 & 1.26 & 0.96\\
fastsae & Beta-ICAR, fastsae (INLA) & 0.0495 & 19.1 & 15.9 & 93.2 (0.1) & 0.193 & 1.25 & 9.48\\
tipsae & Beta-ICAR, tipsae (Stan) & 0.0490 & 18.8 & 15.6 & 94.1 (0.1) & 0.203 & 1.38 & 126\\
\bottomrule
\end{tabular}
\end{table}

All model-based estimators reduced the RMSE of the direct estimator from 0.0725 to between 0.0490 and 0.0518 (Table \ref{tab:mc}). The two beta models were about equally accurate and covered 93.2\% and 94.1\% of the true proportions. The Fay-Herriot EBLUPs covered 94.4\% (spatial: 94.2\%), and their analytical MSE estimates were 1.27 and 1.26 times the empirical MSE in the median area. The largest \(\hat{R}\) of the \CRANpkg{tipsae} fits exceeded 1.01 in 167 and 1.05 in 20 of the 200 replicates. An earlier version of this simulation, with a PC prior on the area standard deviation and the beta precision evaluated at the direct estimates, covered 89.0\% with \texttt{hb\_area()}.

\section{Example: poverty in Emilia-Romagna}\label{sec:applications}

We illustrate a complete workflow with the head-count poverty ratios of the 38 health districts of the Italian region Emilia-Romagna in 2016, distributed with \CRANpkg{tipsae} \citep{denicolo2024tipsae}. The data contain the direct estimates \texttt{hcr}, their sampling variances \texttt{vars}, the sample sizes \texttt{n} (10 to 285 households), the populations \texttt{pop} and the province \texttt{prov} of each district, and a shapefile of the districts. The only covariate, \texttt{x}, is artificial, so the example illustrates the software rather than a substantive analysis.

\begin{verbatim}
library("fastsae")
data("emilia_cs", "emilia_shp", package = "tipsae")
shp <- sf::st_as_sf(emilia_shp)
shp <- shp[match(emilia_cs$id, shp$NAME_DISTRICT), ]
W <- spdep::nb2mat(spdep::poly2nb(shp), style = "B")
dimnames(W) <- list(emilia_cs$id, emilia_cs$id)
\end{verbatim}

We fit a Fay-Herriot model and a beta model with a BYM2 area effect, both with the sampling variances of the direct estimates:

\begin{verbatim}
fit_fh <- eblup_fh(hcr ~ x, vardir = "vars", domain = "id",
  data = emilia_cs, print_result = FALSE)
fit_beta <- hb_area(hcr ~ x, data = emilia_cs, domain = "id",
  vardir = "vars", family = "beta", spatial = "bym2", W = W,
  print_result = FALSE)
\end{verbatim}

\texttt{compare\_sae()} summarizes how far the two sets of estimates agree and which of them has the smaller estimated uncertainty; below are its agreement and efficiency measures:

\begin{verbatim}
cmp <- compare_sae(fit_fh, fit_beta, names = c("FH", "Beta-BYM2"))
subset(cmp$metrics, Group %in% c("Agreement", "Efficiency"), c(Metric, Value))
\end{verbatim}

\begin{verbatim}
#>                                                    Metric             Value
#> 2                        Mean Difference (Beta-BYM2 - FH)            0.0013
#> 3                                 95% Limits of Agreement [-0.0110, 0.0137]
#> 4                          Mean Absolute Difference (MAE)            0.0042
#> 5                     Root Mean Squared Difference (RMSD)            0.0064
#> 6                     Lin's Concordance Correlation (CCC)            0.9604
#> 13                        Mean MSE Ratio (FH / Beta-BYM2)            1.1487
#> 14                      Median MSE Ratio (FH / Beta-BYM2)            1.1611
#> 15 % of domains where Beta-BYM2 reports lower uncertainty             63.2%
#> 16                                        Mean RSE % (FH)            16.58%
#> 17                                 Mean RSE % (Beta-BYM2)            15.35%
\end{verbatim}

The regional poverty ratio, estimated directly from the regional sample, is precise enough to publish, so we benchmark the district estimates to it, and to the provincial direct estimates, with the logit adjustment that keeps them in \((0, 1)\). For the illustration we use population-weighted means of the district direct estimates as targets:

\begin{verbatim}
prov_target <- tapply(emilia_cs$hcr * emilia_cs$pop, emilia_cs$prov, sum) /
  tapply(emilia_cs$pop, emilia_cs$prov, sum)
reg_target <- sum(emilia_cs$hcr * emilia_cs$pop) / sum(emilia_cs$pop)
bm <- benchmark_sae(fit_beta, target = prov_target, group = emilia_cs$prov,
  weight = emilia_cs$pop, national_target = reg_target, method = "logit")
head(bm[, c("domain", "group", "original", "benchmarked")], 4)
\end{verbatim}

\begin{verbatim}
#>                domain group   original benchmarked
#> 1 CASALECCHIO DI RENO    BO 0.05724921  0.05642148
#> 2   CITTA' DI BOLOGNA    BO 0.06635247  0.06540226
#> 3               IMOLA    BO 0.07000428  0.06900564
#> 4         PIANURA EST    BO 0.06648950  0.06553746
\end{verbatim}

Finally, the estimates are classified by their RSE, mapped (Figure \ref{fig:map}), and exported to a spreadsheet with \texttt{export\_sae()}:

\begin{verbatim}
as.data.frame(table(tier = add_reliability_flags(fit_beta)$reliability))
\end{verbatim}

\begin{verbatim}
#>                                 tier Freq
#> 1               Reliable (RSE <= 25%)   38
#> 2 Use with caution (25% < RSE <= 50%)    0
#> 3             Unreliable (RSE > 50%)    0
\end{verbatim}

\begin{verbatim}
p_est <- map_sae(fit_beta, sf_geom = shp, key = "NAME_DISTRICT",
  type = "estimate", title = "Head-count ratio", subtitle = "")
p_rse <- map_sae(fit_beta, sf_geom = shp, key = "NAME_DISTRICT",
  type = "rse", title = "RSE (%)", subtitle = "")
patchwork::wrap_plots(p_est, p_rse)
\end{verbatim}

\begin{figure}

{\centering \includegraphics[width=1\linewidth,alt={Two choropleth maps of the 38 health districts of Emilia-Romagna. Left: estimated poverty ratios. Right: relative standard errors in percent.}]{fastsae_files/figure-latex/map-1} 

}

\caption{Head-count poverty ratios of the health districts of Emilia-Romagna in 2016 estimated by the beta model with a BYM2 area effect (left) and their relative standard errors (right).}\label{fig:map}
\end{figure}

\section{Summary and discussion}\label{sec:summary}

\CRANpkg{fastsae} brings the frequentist and hierarchical Bayes small area models used in official statistics, and the steps that follow a fit, together in one interface. On the same data, the compiled frequentist implementations reproduce the estimates of \CRANpkg{sae} and fit the spatial Fay-Herriot model with analytical MSE for 1,000 areas in 0.52 s, against 24.10 s for \CRANpkg{sae}. The HB functions fit beta and other non-Gaussian models in seconds with INLA, and the same models can be checked against full MCMC with the JAX engine without rewriting them.

Some limitations should be kept in mind. INLA has a fixed cost per fit, so for small Gaussian models numerical integration as in \CRANpkg{hbsae} is faster: for the Fay-Herriot model with 250 areas, \texttt{hb\_area()} took 0.70 s and \texttt{hbsae::fSAE.Area()} 0.05 s (script \texttt{benchmarks/bench\_hb\_normal.R}). The beta and negative binomial dispersions are fixed at first-stage posterior means, and their uncertainty is not propagated; the gamma shape is fixed at the direct estimates. Variance estimates of zero are reported but not prevented, as they would be by adjusted likelihoods \citep{yoshimori2014new}. INLA is an approximation whose accuracy can deteriorate for non-Gaussian likelihoods with very little information per area; the JAX engine allows such cases to be checked. Our benchmarks were run on a single machine.

Future work will add unit-level models for binary and count outcomes that use unit-level population covariates, as the empirical best predictor of \citet{molina2010small} does (the present \texttt{hb\_unit()} predicts the non-sampled units at the mean covariates, which is approximate for these families), a two-stage shape for the gamma model, zero-inflated beta models for indicators with many zeros, and second-order MSE estimators for benchmarked EBLUPs.

\section*{Computational details}\label{computational-details}
\addcontentsline{toc}{section}{Computational details}

The results were obtained with R 4.6.1, \CRANpkg{fastsae} 1.0.0, \pkg{R-INLA} 26.8.7, \CRANpkg{tipsae} 1.0.3, \CRANpkg{sae} 1.3, \CRANpkg{emdi} 2.2.3 and \CRANpkg{hbsae} 1.2; the JAX engine used JAX 0.5.0 and NumPyro 0.19.0. R and the packages are available from CRAN at \url{https://CRAN.R-project.org/}, \pkg{R-INLA} from \url{https://www.r-inla.org/}, and the source of \CRANpkg{fastsae}, including the benchmark scripts, from \url{https://github.com/ridsonap/fastsae}.

\section*{Acknowledgments}\label{acknowledgments}
\addcontentsline{toc}{section}{Acknowledgments}

The authors thank the Department of Computational Statistics and the Department of Statistics at Politeknik Statistika STIS for institutional support, and colleagues at BPS-Statistics Indonesia for their feedback.

\bibliography{fastsae.bib}

\address{%
Ridson Al Farizal P.\\
Department of Computational Statistics, Politeknik Statistika STIS\\%
Jl. Otto Iskandardinata No.~64C\\ Jakarta Timur 13330, Indonesia\\
%
%
\textit{ORCiD: \href{https://orcid.org/0000-0003-0617-0214}{0000-0003-0617-0214}}\\%
\href{mailto:alfrzlp@gmail.com}{\nolinkurl{alfrzlp@gmail.com}}%
}

\address{%
Azka Ubaidillah\\
Department of Statistics, Politeknik Statistika STIS\\%
Jl. Otto Iskandardinata No.~64C\\ Jakarta Timur 13330, Indonesia\\
%
%
\textit{ORCiD: \href{https://orcid.org/0000-0002-3597-0459}{0000-0002-3597-0459}}\\%
\href{mailto:azka@stis.ac.id}{\nolinkurl{azka@stis.ac.id}}%
}
